\ncdmodultestmthree{4}{Familie, By og Fremtid}{全真模考卷四：城乡搬迁、客运投诉与家庭分工}{城乡搬迁选词完形 / 周末改期与子女病假双聊天 / 首次企业实习补句 / 三人谈家务与家庭分工 / 客运大巴投诉与丹麦适应回信 / 家庭分工与丹麦语学习 Mindmap 及假期旅游四图双人互动}{B1 全核心语法总复习：-s 被动与 blive 被动辨析 · 介词 + 形式宾语 det + at/hv- 从句 · 分词作定语 · 无 hvis 条件倒装 · 让步对比 til trods for / selvom / mens... derimod · 篇章连接词 For / Men / Derfor / Desuden}

% =====================================================================
% Del 1: Læsning (50 minutter · 24 point)
% =====================================================================
\ncdsubsec{Del 1: Læsning（阅读理解 · 50 分钟 · 24 分）}

\begin{tipbox}{考场规则与配速提醒（Læsning Opgave 1, 2A, 2B, 3 og 4）}
Opgave 1, 2A, 2B, 3 og 4 udleveres samtidig. Samlet tid: \textbf{50 minutter}. Du må \textbf{ikke} bruge ordbog, noter eller mobiltelefon. Der gives 1 point for hvert korrekt svar; maksimum er \textbf{24 point}（及格线：\textbf{15--24 point = Bestået}；\textbf{0--14 point = Ikke bestået}）。建议时间分配：Opgave 1 (8 min) $\rightarrow$ Opgave 2A \& 2B (10 min) $\rightarrow$ Opgave 3 (14 min) $\rightarrow$ Opgave 4 (13 min)，最后留 5 分钟复查核对。
\end{tipbox}

% ---------------------------------------------------------------------
\subsubsection*{Opgave 1 (6 point): \textit{Fra storbyen til en lille landsby}}

\begin{emailbox}{Kursistark --- Instruktion（Læsning, Opgave 1）}
\begin{itemize}
  \item Læs teksten \textit{Fra storbyen til en lille landsby}.
  \item Skriv de ord (1--6), der mangler. Du skal kun bruge hvert ord \textbf{én gang}.
  \item Der er \textbf{seks ord, du ikke skal bruge}. Se eksemplet (0).
\end{itemize}

\vspace{0.35em}
\hrule
\vspace{0.45em}

\textbf{\color{NCDNavy}Fra storbyen til en lille landsby}

\vspace{0.3em}
Efter syv år i en dyr toværelses lejlighed midt i København besluttede Signe og Kasper, at de ville flytte ud på landet. De havde længe drømt om mere plads og en have, men det var svært at finde et hus til en \textbf{(0)} \underline{\textbf{rimelig}} pris i hovedstaden.

\vspace{0.35em}
En weekend i foråret besøgte de nogle venner i en lille landsby på Fyn, \textbf{(1)} \rule{2.0cm}{0.4pt} der kun bor omkring fire hundrede mennesker. De faldt straks for den rolige stemning og de grønne marker omkring byen. Selvom de var lidt bekymrede for den lange transporttid til arbejde, købte de et gammelt murstenshus med frugttræer i haven.

\vspace{0.35em}
Den første måned var Signe lidt usikker på, \textbf{(2)} \rule{2.0cm}{0.4pt} hun ville savne caféerne og kulturlivet i København. Om aftenen var gaderne helt stille, og i haven kunne de høre fuglene synge, \textbf{(3)} \rule{2.0cm}{0.4pt} de spiste morgenmad sammen. Signe havde egentlig troet, at folk på landet holdt sig lidt for sig selv, \textbf{(4)} \rule{2.0cm}{0.4pt} allerede den første weekend bankede naboerne på døren med friskbagt brød.

\vspace{0.35em}
I dag pendler Kasper med tog fire dage om ugen, og han har endelig vænnet sig \textbf{(5)} \rule{2.0cm}{0.4pt} den ekstra rejsetid, fordi han kan arbejde hjemmefra hver fredag. Parret er stadig meget glade for deres beslutning, \textbf{(6)} \rule{2.0cm}{0.4pt} turen til kontoret faktisk tager længere tid end før. Livet på landet har givet dem meget mere ro i hverdagen.

\vspace{0.55em}
\begin{center}
\fbox{\parbox{0.93\linewidth}{\centering\vspace{4pt}
\textbf{til} \quad\textbar\quad \textbf{hvis} \quad\textbar\quad \textbf{hvor} \quad\textbar\quad \textbf{ikke} \quad\textbar\quad \textbf{mens} \quad\textbar\quad \textbf{rimelig} \quad\textbar\quad \textbf{under}\\[3pt]
\textbf{om} \quad\textbar\quad \textbf{fordi} \quad\textbar\quad \textbf{selvom} \quad\textbar\quad \textbf{lige} \quad\textbar\quad \textbf{men} \quad\textbar\quad \textbf{så}%
\vspace{4pt}}}
\end{center}
\end{emailbox}

% ---------------------------------------------------------------------
\subsubsection*{Opgave 2A (NY · 3 point): \textit{Mette skriver til sin veninde Louise}}

\begin{emailbox}{Kursistark --- Instruktion（Læsning, Opgave 2A）}
\textbf{Situation:} Mette skriver til sin veninde Louise.
\begin{itemize}
  \item Læs Mettes og Louises chat. Der er 5 beskeder fra Mette.
  \item Find de svar fra Louise, der passer (A--F).
  \item Skriv bogstavet på det svar, der passer. Der er \textbf{tre svar, du ikke skal bruge}. Se eksemplet (0).
\end{itemize}
\end{emailbox}

\vspace{0.25em}
{\small
\renewcommand{\arraystretch}{1.35}
\begin{longtable}{@{} >{\raggedright\arraybackslash}p{8.25cm} @{\hspace{0.5cm}} >{\raggedright\arraybackslash}p{8.25cm} @{}}
\toprule
\textbf{\color{NCDNavy}Chat mellem Mette og Louise} & \textbf{\color{NCDNavy}Louises mulige svar (Z, A--F)} \\
\midrule
\endfirsthead
\toprule
\textbf{\color{NCDNavy}Chat mellem Mette og Louise (fortsat)} & \textbf{\color{NCDNavy}Louises mulige svar (fortsat)} \\
\midrule
\endhead
\bottomrule
\endfoot

\textbf{\color{NCDNavy}Mette:}\newline
Hej Louise! Glæder du dig til vores tur til Aarhus på lørdag med biograf og kunstmuseum?\par\vspace{4pt}
\textbf{0.} \textbf{\color{NCDTeal}Louise:} \textcolor{NCDRed}{\textbf{[ Eksempel: Z ]}} &
\textbf{\color{NCDRed}Z:}\ Ja, helt vildt! Jeg har set frem til det hele ugen. \\
\midrule

\textbf{\color{NCDNavy}Mette:}\newline
Også mig, men jeg har lige set på Rejseplanen, at der er sporarbejde på strækningen hele lørdagen. Så toget kører slet ikke.\par\vspace{4pt}
\textbf{1.} \textbf{\color{NCDTeal}Louise:} \rule{1.2cm}{0.4pt} &
\textbf{\color{NCDNavy}A:}\ Det lyder godt! Jeg elsker nemlig togbussen, for så kan man sove lidt ekstra undervejs.\par\vspace{6pt}
\textbf{\color{NCDNavy}B:}\ Søndag passer mig fint! Men jeg har desværre ikke tid til både filmen og udstillingen, fordi jeg skal besøge min bror sent om eftermiddagen. Skal vi droppe kunstmuseet? \\
\midrule

\textbf{\color{NCDNavy}Mette:}\newline
Det har du desværre ret i --- det tager næsten dobbelt så lang tid med togbus. Hvad hvis vi tager turen om søndagen i stedet?\par\vspace{4pt}
\textbf{2.} \textbf{\color{NCDTeal}Louise:} \rule{1.2cm}{0.4pt} &
\textbf{\color{NCDNavy}C:}\ Desværre, på søndag arbejder jeg hele dagen fra morgen til aften, så vi må vente til næste måned.\par\vspace{6pt}
\textbf{\color{NCDNavy}D:}\ Øv, hvor ærgerligt! Så skal vi vel med togbus, og det tager jo en evighed. \\
\midrule

\textbf{\color{NCDNavy}Mette:}\newline
Nej, så går vi jo glip af den nye udstilling, og det vil jeg nødigt. Kunne vi så ikke se kunstmuseet om formiddagen og droppe biografen?\par\vspace{4pt}
\textbf{3.} \textbf{\color{NCDTeal}Louise:} \rule{1.2cm}{0.4pt} &
\textbf{\color{NCDNavy}E:}\ Jeg tager gerne begge biografbilletter med på lørdag, så vi ikke kommer for sent til filmen.\par\vspace{6pt}
\textbf{\color{NCDNavy}F:}\ Jo, det lyder som en rigtig god plan! Så kan vi også nå en kop kaffe bagefter. Bestiller du billetter til museet? \\
\midrule

\textbf{\color{NCDNavy}Mette:}\newline
Super, så køber jeg museumsbilletter til os til kl. 10.30. Vi ses på stationen søndag morgen! &
\textit{\color{NCDGray}(Husk: Der er tre svar blandt A--F, som du ikke skal bruge.)} \\
\end{longtable}
}

% ---------------------------------------------------------------------
\subsubsection*{Opgave 2B (NY · 3 point): \textit{Anders skriver til sin chef Camilla}}

\begin{emailbox}{Kursistark --- Instruktion（Læsning, Opgave 2B）}
\textbf{Situation:} Anders skriver til sin chef Camilla.
\begin{itemize}
  \item Læs Anders' og Camillas chat. Der er 5 beskeder fra Anders.
  \item Find de svar fra Camilla, der passer (A--F).
  \item Skriv bogstavet på det svar, der passer. Der er \textbf{tre svar, du ikke skal bruge}. Se eksemplet (0).
\end{itemize}
\end{emailbox}

\vspace{0.25em}
{\small
\renewcommand{\arraystretch}{1.35}
\begin{longtable}{@{} >{\raggedright\arraybackslash}p{8.25cm} @{\hspace{0.5cm}} >{\raggedright\arraybackslash}p{8.25cm} @{}}
\toprule
\textbf{\color{NCDNavy}Chat mellem Anders og Camilla} & \textbf{\color{NCDNavy}Camillas mulige svar (Z, A--F)} \\
\midrule
\endfirsthead
\toprule
\textbf{\color{NCDNavy}Chat mellem Anders og Camilla (fortsat)} & \textbf{\color{NCDNavy}Camillas mulige svar (fortsat)} \\
\midrule
\endhead
\bottomrule
\endfoot

\textbf{\color{NCDNavy}Anders:}\newline
Godmorgen Camilla. Min søn Noah er desværre vågnet med høj feber, så jeg bliver nødt til at tage barns første sygedag i dag.\par\vspace{4pt}
\textbf{0.} \textbf{\color{NCDTeal}Camilla:} \textcolor{NCDRed}{\textbf{[ Eksempel: Z ]}} &
\textbf{\color{NCDRed}Z:}\ Godmorgen Anders. Det er jeg ked af at høre --- god bedring til ham! Kan du overhovedet arbejde lidt hjemmefra? \\
\midrule

\textbf{\color{NCDNavy}Anders:}\newline
Ja, min kone kommer hjem ved frokosttid og tager over, så jeg kan godt arbejde hjemmefra i eftermiddag. Men hvad med mit kundemøde kl. 9.30?\par\vspace{4pt}
\textbf{1.} \textbf{\color{NCDTeal}Camilla:} \rule{1.2cm}{0.4pt} &
\textbf{\color{NCDNavy}A:}\ Perfekt! Husk også lige at registrere barns første sygedag i vores system, ikke?\par\vspace{6pt}
\textbf{\color{NCDNavy}B:}\ Nej, Martin har desværre ferie hele denne uge, så du må selv ringe og aflyse mødet med kunden. \\
\midrule

\textbf{\color{NCDNavy}Anders:}\newline
Tusind tak, det er en kæmpe lettelse! Jeg sender ham straks alle noterne og præsentationen på mail.\par\vspace{4pt}
\textbf{2.} \textbf{\color{NCDTeal}Camilla:} \rule{1.2cm}{0.4pt} &
\textbf{\color{NCDNavy}C:}\ Ja, hvis du kan nå at gøre månedsbudgettet færdigt inden fyraften, ville det være en stor hjælp.\par\vspace{6pt}
\textbf{\color{NCDNavy}D:}\ Så synes jeg, du bare skal holde helt fri resten af dagen og først komme på arbejde igen på mandag. \\
\midrule

\textbf{\color{NCDNavy}Anders:}\newline
Ja, jeg har allerede registreret det i vores tidsystem, og jeg logger på Teams fra kl. 13. Er der noget særligt, jeg skal prioritere i eftermiddag?\par\vspace{4pt}
\textbf{3.} \textbf{\color{NCDTeal}Camilla:} \rule{1.2cm}{0.4pt} &
\textbf{\color{NCDNavy}E:}\ Bare rolig, jeg spørger Martin, om han kan tage kundemødet for dig. Kan du sende ham dit materiale?\par\vspace{6pt}
\textbf{\color{NCDNavy}F:}\ Det behøver du ikke, for jeg har allerede ringet til din kone og aftalt, at hun bliver hjemme hele dagen. \\
\midrule

\textbf{\color{NCDNavy}Anders:}\newline
Det skal jeg nok gøre færdigt før kl. 16. Vi snakkes ved på Teams i eftermiddag! &
\textit{\color{NCDGray}(Husk: Der er tre svar blandt A--F, som du ikke skal bruge.)} \\
\end{longtable}
}

% ---------------------------------------------------------------------
\subsubsection*{Opgave 3 (NY · 5 point): \textit{Fatimas første praktikplads}}

\begin{emailbox}{Kursistark --- Instruktion（Læsning, Opgave 3）}
\begin{itemize}
  \item Læs teksten \textit{Fatimas første praktikplads}. Der er fem afsnit.
  \item I hvert afsnit mangler der en sætning. Under hvert afsnit står der fire sætninger.
  \item Skriv bogstavet på den sætning, der passer i afsnittet. Der er \textbf{tre sætninger, du ikke skal bruge} i hvert afsnit.
\end{itemize}

\vspace{0.35em}
\hrule
\vspace{0.45em}

\textbf{1.} Fatima har boet i Odense i to år og har længe drømt om at arbejde med bogholderi og økonomi i Danmark. I sidste måned bestod hun sin danskprøve og fik endelig tilbudt en otte ugers praktikplads på et lokalt regnskabskontor. \rule{3.2cm}{0.4pt}. For selvom hun har arbejdet på kontor i sit hjemland, er alle danske regler og ord helt nye for hende.
\begin{enumerate}[label=\textbf{\Alph*.}, leftmargin=1.8em, itemsep=0.12em, topsep=0.25em]
  \item Hun er både meget glad og lidt nervøs før sin første dag.
  \item Hun har derfor besluttet at sige nej tak til praktikpladsen.
  \item Derfor kender hun allerede alle de danske skatteregler rigtig godt.
  \item Men hendes chef vil helst ikke have praktikanter på kontoret.
\end{enumerate}

\vspace{0.4em}
\textbf{2.} Mandag morgen bliver Fatima budt velkommen af sin kontaktperson, Mette, som viser hende rundt på kontoret og præsenterer hende for de andre ansatte. De første par dage skal Fatima læse gamle fakturaer og lære virksomhedens digitale system at kende. \rule{3.2cm}{0.4pt}. Heldigvis siger Mette, at hun altid må spørge, hvis der er et fagord eller en opgave, hun ikke forstår.
\begin{enumerate}[label=\textbf{\Alph*.}, leftmargin=1.8em, itemsep=0.12em, topsep=0.25em]
  \item Det er helt nemt for hende, fordi hun bruger systemet hjemmefra hver dag.
  \item Det tager lidt længere tid, end hun havde regnet med, fordi mange ord er svære.
  \item Derfor beder Mette hende om at arbejde helt alene uden at tale med nogen.
  \item Så Fatima bliver sendt hjem allerede før frokost den første dag.
\end{enumerate}

\vspace{0.4em}
\textbf{3.} Efter to uger har Fatima fået mere styr på computerprogrammet, men pauserne i kantinen føles stadig lidt svære. Når de otte kollegaer sidder sammen og spiser frokost, taler de nemlig hurtigt og griner af indforståede jokes. \rule{3.2cm}{0.4pt}. Derfor beslutter hun sig for at sætte sig ved siden af en venlig kollega og selv stille et par spørgsmål om weekendens planer.
\begin{enumerate}[label=\textbf{\Alph*.}, leftmargin=1.8em, itemsep=0.12em, topsep=0.25em]
  \item Fatima foretrækker faktisk at spise sin frokost alene ved sit skrivebord hver dag.
  \item Kollegaerne taler kun engelsk i kantinen, så Fatima kan forstå det hele.
  \item Det gør hende nogle gange usikker på, hvornår hun selv skal sige noget.
  \item Men det er kun, fordi hun ikke kan lide maden i kantinen.
\end{enumerate}

\vspace{0.4em}
\textbf{4.} Midt i praktikperioden får Fatima lov til selv at kontrollere månedens kvitteringer og taste tal ind i regnskabet. En eftermiddag opdager hun pludselig, at der mangler et stort beløb i en af mapperne. \rule{3.2cm}{0.4pt}. Mette bliver meget imponeret over, at Fatima har opdaget fejlen, og roser hende for at være så omhyggelig.
\begin{enumerate}[label=\textbf{\Alph*.}, leftmargin=1.8em, itemsep=0.12em, topsep=0.25em]
  \item Hun tænker, at det nok ikke er så vigtigt, og lægger bare mappen væk.
  \item Fatima bliver bange for, at hun selv har ødelagt det hele, og siger ingenting.
  \item Efter en halv times søgen finder hun heldigvis det forsvundne bilag og viser det til Mette.
  \item Derfor ringer hun straks til en anden virksomhed for at søge et nyt job.
\end{enumerate}

\vspace{0.4em}
\textbf{5.} Da de otte uger næsten er gået, bliver Fatima kaldt ind til en samtale på chefens kontor. Hun er spændt på, hvad chefen vil sige om hendes indsats i løbet af praktikken. \rule{3.2cm}{0.4pt}. Fatima går hjem med et kæmpe smil, for nu ved hun, at drømmen om et fast arbejde i Danmark faktisk kan blive til virkelighed.
\begin{enumerate}[label=\textbf{\Alph*.}, leftmargin=1.8em, itemsep=0.12em, topsep=0.25em]
  \item Chefen fortæller hende desværre, at kontoret ikke kan bruge hendes hjælp længere.
  \item Chefen tilbyder hende et vikariat på deltid fra næste måned, fordi teamet er så glade for hende.
  \item De bliver enige om, at Fatima slet ikke passer til at arbejde med regnskab.
  \item Men chefen har glemt, at Fatima overhovedet har været i praktik hos dem.
\end{enumerate}
\end{emailbox}

% ---------------------------------------------------------------------
\subsubsection*{Opgave 4 (NY · 7 point): \textit{Tre personer fortæller om arbejdsdeling i hjemmet}}

\begin{emailbox}{Kursistark --- Instruktion \& Svarskema（Læsning, Opgave 4）}
\begin{itemize}
  \item Læs de tre tekster, og læs spørgsmålene (1--7).
  \item Find den person, der passer til hvert af de syv spørgsmål.
  \item Sæt kryds (\textbf{X}) ud for den person (A, B eller C), der passer til spørgsmålet. Se eksemplet (0).
\end{itemize}

\vspace{0.3em}
{\small
\renewcommand{\arraystretch}{1.3}
\begin{tabularx}{\linewidth}{@{} >{\bfseries}p{0.6cm} >{\raggedright\arraybackslash}X >{\centering\arraybackslash}p{1.65cm} >{\centering\arraybackslash}p{1.65cm} >{\centering\arraybackslash}p{1.65cm} @{}}
\toprule
 & \textbf{Spørgsmål} & \textbf{Stine} & \textbf{Karim} & \textbf{Elena} \\
\midrule
0. & Hvem mener, at det er godt at fordele opgaverne efter evner og interesser? & A \underline{\hspace{0.45cm}} & B \underline{\hspace{0.45cm}} & C \underline{\textbf{X}} \\
\midrule
1. & Hvem får professionel hjælp udefra til rengøringen? & A \underline{\hspace{0.45cm}} & B \underline{\hspace{0.45cm}} & C \underline{\hspace{0.45cm}} \\
2. & Hvem bruger en fast skriftlig plan i køkkenet til at fordele pligterne? & A \underline{\hspace{0.45cm}} & B \underline{\hspace{0.45cm}} & C \underline{\hspace{0.45cm}} \\
3. & Hvem har overtaget al madlavningen, fordi partneren ikke var god til det? & A \underline{\hspace{0.45cm}} & B \underline{\hspace{0.45cm}} & C \underline{\hspace{0.45cm}} \\
4. & Hvem må ofte bede sine børn flere gange om at gøre deres pligter? & A \underline{\hspace{0.45cm}} & B \underline{\hspace{0.45cm}} & C \underline{\hspace{0.45cm}} \\
5. & Hvem synes, det er irriterende, at hjemmearbejde fører til ekstra husarbejde? & A \underline{\hspace{0.45cm}} & B \underline{\hspace{0.45cm}} & C \underline{\hspace{0.45cm}} \\
6. & Hvem køber ind til familien, mens partneren laver det meste af aftensmaden? & A \underline{\hspace{0.45cm}} & B \underline{\hspace{0.45cm}} & C \underline{\hspace{0.45cm}} \\
7. & Hvem overtog madlavningen i hverdagene, fordi partneren fik skiftende arbejdstider? & A \underline{\hspace{0.45cm}} & B \underline{\hspace{0.45cm}} & C \underline{\hspace{0.45cm}} \\
\bottomrule
\end{tabularx}
}
\end{emailbox}

\begin{emailbox}{Tre personer fortæller om arbejdsdeling i hjemmet (A--C)}
\textbf{\color{NCDNavy}A. Stine}\\
Da vores to børn var helt små, arbejdede jeg på deltid, så jeg lavede næsten alt husarbejde selv. Det førte dog til mange skænderier mellem min mand Nikolaj og mig, fordi jeg altid var træt om aftenen. I dag arbejder vi begge to fuldtid, og derfor har vi lavet et fast ugeskema på køleskabet. Nikolaj støvsuger og gør badeværelset rent, mens jeg tager mig af tøjvasken og indkøbene. Til gengæld elsker Nikolaj at stå i køkkenet, så han laver aftensmad fem dage om ugen, og jeg laver bare noget hurtigt i weekenden. Vores børn på seks og ni år har også faste pligter: De skal selv rydde op på deres værelser og tømme opvaskemaskinen et par gange om ugen. Jeg synes ikke, at ligestilling i hjemmet handler om at tælle hvert eneste minut, men om at begge forældre får lov til at slappe af om aftenen.

\vspace{0.45em}
\textbf{\color{NCDNavy}B. Karim}\\
I min barndom lavede min mor al mad og rengøring, så da jeg blev gift med Sarah, regnede jeg ærligt talt med, at det ville blive på samme måde. Det ændrede sig hurtigt, da Sarah begyndte som sygeplejerske med skiftende nattevagter. I dag handler jeg næsten altid ind og laver varm aftensmad alle hverdage, fordi hun ofte er helt udkørt efter arbejde. Ingen af os har lyst til at bruge vores frie weekender på at skrubbe gulve, så vi har valgt at betale et rengøringsfirma, der kommer hver anden fredag. Vores tre teenagere gider desværre ikke at lave mad, men til gengæld skiftes de til at slå græs, bære skrald ud og lægge tøj sammen --- selvom jeg stadig må minde dem om det flere gange. Efter min mening skal mænd tage lige så meget ansvar i familien som kvinder.

\vspace{0.45em}
\textbf{\color{NCDNavy}C. Elena}\\
Jeg bor sammen med min danske kæreste Frederik, og vi har ingen børn endnu. Da jeg arbejder hjemmefra tre dage om ugen, kommer jeg nogle gange automatisk til at sætte en vask over i frokostpausen, mens Frederik pendler to timer hver dag. Det irriterer mig faktisk lidt, for hjemmearbejde er jo stadig rigtigt arbejde! Tidligere prøvede vi at skiftes til at lave mad, men Frederik brændte maden på to gange og hader at følge opskrifter. Nu står jeg for al madlavningen, fordi jeg godt kan lide det. Til gengæld tager Frederik altid opvasken bagefter, pudser alle vinduer og kører affald på genbrugspladsen. Jeg tror, at god ligestilling handler om at dele opgaverne efter, hvad man hver især er bedst til --- ikke at begge skal lave præcis de samme ting. Hvis vi får børn en dag, skal de tidligt lære at hjælpe til.
\end{emailbox}

% =====================================================================
% Del 2: Skrivning (双任务：Opgave 1 投诉信 + Opgave 2 私人回信 >= 90 ord)
% =====================================================================
\clearpage
\ncdsubsec{Del 2: Skrivning（书面表达 · 双任务实战）}

\subsubsection*{Opgave 1（正式/半正式应用文 · 投诉信 Klage）}

\begin{emailbox}{Kursistark --- Skrivning, Opgave 1}
\textbf{Situation:} Du tog i går med langdistancebussen \textit{NordBus} fra København til Aalborg, men rejsen var meget dårlig. Bussen var kraftigt forsinket uden besked, dit sæde var ødelagt, og der var intet fungerende toilet. Da du klagede, var servicen også dårlig.\\
\textbf{Opgave:} Skriv en klage til \textit{NordBus}. Du skal fortælle:
\begin{itemize}[itemsep=0.05em, topsep=0.08em, parsep=0pt]
  \item hvor og hvornår du rejste
  \item hvad der gik galt under turen
  \item hvordan chaufføren eller kundeservice reagerede
  \item hvordan \textit{NordBus} kan kompensere dig.
\end{itemize}

\vspace{0.15em}
\hrule
\vspace{0.15em}
\textbf{Klage til NordBus}\\[2pt]
\rule{\linewidth}{0.4pt}\\[2pt]
\rule{\linewidth}{0.4pt}\\[2pt]
\rule{\linewidth}{0.4pt}
\end{emailbox}

\subsubsection*{Opgave 2（私人电子邮件回复 · Minimum 90 ord · Tid: 30 minutter）}

\begin{emailbox}{Kursistark --- Skrivning, Opgave 2}
\textbf{Situation:} Du har fået e-mailen her nedenfor fra din veninde Line. Læs e-mailen, og skriv et svar.\\
\textbf{Opgave:} Skriv e-mailen til Line. Du skal svare på alle Lines spørgsmål og indlede og afslutte e-mailen på en passende måde.\\
\textbf{Længde:} Minimum 90 ord. \qquad \textbf{Tid:} 30 minutter.

\vspace{0.2em}
\hrule
\vspace{0.2em}
\hfill\textit{Odense d. 19. november 2025}\\[2pt]
\textit{Kære \ldots}

\vspace{0.15em}
\textit{Hvordan har du det? Vi har ikke skrevet sammen i lang tid, og jeg tænkte på, hvordan det går med dit liv i Danmark nu sammenlignet med i begyndelsen. Hvad overrasker dig egentlig mest i Danmark? Hvad er blevet nemmere for dig i hverdagen? Er der stadig noget ved dansk kultur eller det danske sprog, som du synes er svært? Og hvilke planer har du for de næste måneder?}

\vspace{0.15em}
\textit{Jeg glæder mig rigtig meget til at høre fra dig!}

\vspace{0.15em}
\textit{Kærlig hilsen\\
Line}

\vspace{0.15em}
\hrule
\vspace{0.15em}
\textbf{Til:} Line \qquad \textbf{Emne:} Svar: Livet i Danmark nu\\[2pt]
\rule{\linewidth}{0.4pt}\\[2pt]
\rule{\linewidth}{0.4pt}\\[2pt]
\rule{\linewidth}{0.4pt}
\end{emailbox}

% =====================================================================
% Del 3: Mundtlig kommunikation
% =====================================================================
\clearpage
\ncdsubsec{Del 3: Mundtlig kommunikation（口语交际全真题卡）}

\subsubsection*{Mundtlig Opgave 1: \textit{Emneopgave med Mindmap}}

\begin{tipbox}{官方考场规则与红线提醒（Mundtlig Opgave 1）}
考生须自带 \textbf{2 张} 已填好的 Mindmap（须选自工作/实习/学习/语言学校或市民社会话题）。\textbf{考场铁律：Mindmap 上只能写关键词或短语（stikord），绝不能写完整句子！} 考官看你的空白/关键词卡记笔记：Del 1 脱稿陈述 \textbf{1--2 分钟}，Del 2 围绕你的内容追问 \textbf{3--4 分钟}（须简要说明理由，如 \textit{fordi / derfor / synes}）。\\
\textbf{\color{NCDRed}官方禁选红线：}\textit{Min bolig/lejlighed}、\textit{Min familie} 与 \textit{Min hverdag} 因复杂度不足 \textbf{禁止用于 Modul 3 Opgave 1}！下面的 \textit{Arbejdsdeling i min familie} 聚焦分工、子女责任与家庭平权，为 SIRI 官方推荐 B1 议题。
\end{tipbox}

\begin{grammatikbox}{Mindmap 4A：Arbejdsdeling i min familie（官方 6 辐射格 · 纯 Stikord 备考卡）}
\small
\renewcommand{\arraystretch}{1.32}
\begin{tabularx}{\linewidth}{@{} >{\bfseries\color{NCDNavy}\raggedright\arraybackslash}p{4.3cm} >{\raggedright\arraybackslash}X @{}}
\toprule
\textbf{Mindmap 官方分支} & \textbf{Stikord 关键词笔记（禁止整句！）} \\
\midrule
1. hvem bor sammen & mig, min partner (Thomas), datter (8 år) · begge fuldtidsarbejde \\
2. madlavning og indkøb & skiftes hver anden dag · ugentlig madplan · fælles indkøb lørdag \\
3. rengøring og tøjvask & fast lørdagsrutine · støvsuge, badeværelse · vaske og lægge tøj sammen \\
4. børns pligter/ansvar & rydde op på værelse · dække bord · tømme opvaskemaskine \\
5. økonomi/planlægning & fælles budgetkonto · dele faste udgifter · familiekalender på mobilen \\
6. mening om ligestilling i hjemmet & vigtigt med 50/50 · begge har fritid · godt forbillede for børn \\
\bottomrule
\end{tabularx}

\vspace{0.45em}
\textbf{\color{NCDRed}Del 2 考官追问预测（Eksaminators opfølgende spørgsmål）：}
\begin{enumerate}
  \item \textit{Du siger, at I skiftes til at lave mad hver anden dag. Hvad laver du helst, når det er din tur?}
  \item \textit{Hvordan reagerer jeres datter på sine pligter --- gør hun dem altid uden problemer?}
  \item \textit{Er der nogle huslige opgaver, som du eller din partner slet ikke kan lide at lave?}
  \item \textit{Hvordan har arbejdsdelingen i dit hjem ændret sig i forhold til, hvordan du selv voksede op?}
\end{enumerate}
\end{grammatikbox}

\begin{grammatikbox}{Mindmap 4B：Mine erfaringer med at lære dansk（官方 6 辐射格 · 纯 Stikord 备考卡）}
\small
\renewcommand{\arraystretch}{1.32}
\begin{tabularx}{\linewidth}{@{} >{\bfseries\color{NCDNavy}\raggedright\arraybackslash}p{4.3cm} >{\raggedright\arraybackslash}X @{}}
\toprule
\textbf{Mindmap 官方分支} & \textbf{Stikord 关键词笔记（禁止整句！）} \\
\midrule
1. sprogskole/modul & aftenhold to gange ugentligt · DU3 Modul 3 · 16 kursister i klassen \\
2. tid og metoder & 45 min. lektier dagligt · podcast i toget · skrive nye gloser op \\
3. udtale og grammatik & svære bløde d'er og vokaler · øve ordstilling og passiv \\
4. dansk i hverdagen/job & snakke med kollegaer i frokostpausen · bestille i butikker \\
5. største udfordring & forstå hurtig tale · genert over fejl · skifte fra engelsk \\
6. mening om udvikling & meget mere selvsikker nu · stolt over fremskridt · mål: PD3 \\
\bottomrule
\end{tabularx}

\vspace{0.45em}
\textbf{\color{NCDRed}Del 2 考官追问预测（Eksaminators opfølgende spørgsmål）：}
\begin{enumerate}
  \item \textit{Du nævner, at du lytter til podcasts i toget. Hvilke podcasts eller medier kan du anbefale?}
  \item \textit{Hvad gør du, når en dansker pludselig skifter til engelsk midt i en samtale?}
  \item \textit{Hvad synes du har været det vigtigste, I har lært på Modul 3 indtil nu?}
  \item \textit{Hvordan vil du fortsætte med at bruge og forbedre dit danske efter modultesten?}
\end{enumerate}
\end{grammatikbox}

\clearpage
\subsubsection*{Mundtlig Opgave 2: \textit{Paropgave \& Individuel samtale --- EMNE: Ferie og rejser}}

\begin{emailbox}{Kursistark --- Opgaveark: EMNE: Ferie og rejser（四图互动题卡）}
\small
\renewcommand{\arraystretch}{1.35}
\begin{tabularx}{\linewidth}{@{} >{\raggedright\arraybackslash}X @{\hspace{0.45cm}} >{\raggedright\arraybackslash}X @{}}
\toprule
\textbf{\color{NCDNavy}Billede 1} & \textbf{\color{NCDNavy}Billede 2} \\
\textbf{[图片 1: 丹麦海边度假屋 Sommerhus]}\newline
\textit{En familie eller vennegruppe ved et hyggeligt dansk sommerhus tæt på kysten og sandstranden.} &
\textbf{[图片 2: 欧洲大城市文化观光]}\newline
\textit{Storbyferie i en europæisk hovedstad med museer, caféer, historiske gader og sightseeing.} \\[6pt]
\midrule
\textbf{\color{NCDNavy}Billede 3} & \textbf{\color{NCDNavy}Billede 4} \\
\textbf{[图片 3: 山林自然徒步露营]}\newline
\textit{Aktiv naturferie med vandrestøvler, rygsæk, telt eller træhytte i bjerge og skove.} &
\textbf{[图片 4: 居家休息 Staycation / 一日短途游]}\newline
\textit{Ferie hjemme med afslapning på altanen eller i haven samt små dagsture i lokalområdet med tog eller cykel.} \\
\bottomrule
\end{tabularx}
\end{emailbox}

\begin{grammatikbox}{Testerark --- 官方考官引导脚本与个人深挖三问（Opgave 2）}
\textbf{\color{NCDNavy}Del 1 (Paropgave --- ca. 4 min.)}

Opgavearket lægges foran kursisterne, der præsenteres for emnet på følgende måde:
\begin{quote}\itshape\small
Emnet er ”Ferie og rejser”. Her er nogle billeder, der viser forskellige måder at holde ferie på. I skal tale med hinanden om, hvad I synes om de forskellige måder at holde ferie på, og hvorfor. I skal stille spørgsmål til hinanden og svare på hinandens spørgsmål.
\end{quote}

Hvis kursisterne ikke udtrykker deres præferencer i forhold til \textbf{alle fire billeder}, sørger tester for, at dette sker ved at spørge:
\begin{quote}\itshape\small
Kan I lide at holde ferie i et sommerhus i Danmark / tage på storbyferie i udlandet / tage på aktiv naturferie med vandring eller camping / holde ferie hjemme med korte dagsture? Hvorfor? / Hvorfor ikke?
\end{quote}

\vspace{0.35em}
\hrule
\vspace{0.35em}

\textbf{\color{NCDRed}Del 2 (Individuel opgave --- ca. 3 min. pr. kursist)}

Tester stiller hver kursist nedenstående spørgsmål:
\begin{enumerate}
  \item \textit{Hvordan kan du bedst lide at holde ferie --- og hvorfor?}
  \item \textit{Hvad er vigtigst for dig på en ferie (fx pris, vejr, afslapning eller nye oplevelser)?}
  \item \textit{Kan du fortælle lidt om den bedste ferie, du har haft for nylig? (hvor, hvornår, med hvem, hvad lavede I)}
\end{enumerate}
{\small\textit{Hvis kursisten har besvaret et spørgsmål i Del 1, stilles uddybende spørgsmål, fx: ”Du siger, at \ldots{} vil du fortælle lidt mere om \ldots?”}}
\end{grammatikbox}

% =====================================================================
% Del 4: Retteark, Detaljeret Analyse & Højscore-Modeller
% =====================================================================
\clearpage
\ncdsubsec{Del 4: Retteark, Detaljeret Analyse \& Højscore-Modeller}

\subsubsection*{4.1 官方 \textit{Retteark} 评分卡与考前 30 秒终极自查清单}

\begin{grammatikbox}{DU3 Modul 3 --- Modultest 4 官方标准评分卡（Retteark）}
\small
\renewcommand{\arraystretch}{1.24}
\begin{tabularx}{\linewidth}{@{} >{\bfseries\color{NCDNavy}\raggedright\arraybackslash}p{3.45cm} >{\raggedright\arraybackslash}X >{\centering\arraybackslash}p{1.65cm} >{\centering\arraybackslash}p{2.55cm} @{}}
\toprule
\textbf{考核题型} & \textbf{标准答案（Rettenøgle）} & \textbf{满分} & \textbf{及格标准} \\
\midrule
Læsning Opgave 1 & \textbf{1.} hvor \quad \textbf{2.} om \quad \textbf{3.} mens \quad \textbf{4.} men \quad \textbf{5.} til \quad \textbf{6.} selvom\newline（干扰词：fordi, hvis, ikke, lige, så, under） & 6 point & 分项计分 \\
\cmidrule{1-4}
Læsning Opgave 2A & \textbf{1.} D \quad \textbf{2.} B \quad \textbf{3.} F \quad（干扰项：A, C, E） & 3 point & 分项计分 \\
\cmidrule{1-4}
Læsning Opgave 2B & \textbf{1.} E \quad \textbf{2.} A \quad \textbf{3.} C \quad（干扰项：B, D, F） & 3 point & 分项计分 \\
\cmidrule{1-4}
Læsning Opgave 3 & \textbf{1.} A \quad \textbf{2.} B \quad \textbf{3.} C \quad \textbf{4.} C \quad \textbf{5.} B & 5 point & 分项计分 \\
\cmidrule{1-4}
Læsning Opgave 4 & \textbf{1.} B \quad \textbf{2.} A \quad \textbf{3.} C \quad \textbf{4.} B \quad \textbf{5.} C \quad \textbf{6.} A \quad \textbf{7.} B & 7 point & 分项计分 \\
\cmidrule{1-4}
\textbf{\color{NCDRed}Læsning 总计} & \textbf{Opgave 1, 2A, 2B, 3 og 4 合订总评}（50 分钟 · 禁用词典） & \textbf{24 point} & \textbf{\color{NCDRed}15--24 = Bestået} \\
\midrule
Skrivning Opgave 1 & 向 \textit{NordBus} 投诉：覆盖全部 4 要点 + 语体得体（用 \textit{I / jer / jeres}）+ 句法流畅 & 合格制 & 建议目标 105--145 词 \\
\midrule
Skrivning Opgave 2 & 回复 \textit{Line} 邮件：完整回答全部 4 个问题 + 得体称呼与落款 + \textbf{minimum 90 ord}（30 min.） & 合格制 & \textbf{$\ge$ 90 词}（建议 105--140 词） \\
\midrule
Mundtlig Opgave 1 & 2 张纯 \textit{stikord} Mindmap 抽签其一：Del 1 陈述 1--2 min + Del 2 追问 3--4 min & 合格制 & 整体评估 \textit{Bestået} \\
\midrule
Mundtlig Opgave 2 & \textit{Ferie og rejser}：Del 1 双人围绕全部 4 图互动约 4 min + Del 2 个人问答约 3 min & 合格制 & 整体评估 \textit{Bestået} \\
\bottomrule
\end{tabularx}
\end{grammatikbox}

\begin{tipbox}{考前 30 秒终极自查清单（DU3 Modul 3 结业实战总检）}
作为本册四套 DU3 Modul 3 全真模考的收官卷，进考场前请逐条默念以下 5 条结业保命规则：
\begin{enumerate}
  \item \textbf{阅读 50 分钟禁用词典与 15/24 及格底线}：五部分一次发放（Opg 1 6 分 + 2A/2B 各 3 分 + Opg 3 5 分 + Opg 4 7 分），\textbf{15 分即及格}。遇到 Opg 4 千万别靠“原词匹配”秒选——一定要比对\textbf{同义转述（\textit{Parafrasering}）}与否定反转陷阱！
  \item \textbf{Opg 2A/2B 与 Opg 3 语篇衔接密码}：做聊天补全紧盯\textbf{跨回合指代与语气呼应}（如 \textit{Det har du ret i / ham / sådan noget}）；做段内补句必须“前看一句、后看信号词”（\textit{For} 引出前句原因，\textit{Men} 构成逻辑反转，\textit{Derfor/Så} 顺接结果）。
  \item \textbf{两大 B1 语法引擎与被动语态区分}：\ding{172} 一般规则、习惯或公告用 \textbf{\textit{-s} 被动}（\textit{aftaler planlægges}），具体一次性已发生事件用 \textbf{\textit{blive} 被动}（\textit{bussen blev forsinket / passagerer blev ikke informeret}）；\ding{173} 介词搭配 \textit{at/hv-} 从句时自查衔接（\textit{usikker på, om... / regne med, at...}），并确保从句内 \textit{ikke/altid/faktisk} 严格位于\textbf{动词前}。
  \item \textbf{写作双任务硬核底线}：Opgave 2 私人回信\textbf{严格要求 $\ge$ 90 词}（建议 105--135 词），必须包含礼貌称呼、落款并逐一回答朋友的全部问题；Opgave 1 投诉信或推荐信务必不漏 \textbf{4 个指定要点}，对公司机构统一称呼 \textbf{\textit{I / jer / jeres}}。
  \item \textbf{口试两条不可触碰的红线}：Opgave 1 \textbf{严禁}使用 \textit{Min bolig}、\textit{Min familie}、\textit{Min hverdag} 等低复杂度话题，Mindmap 上只写纯 \textbf{\textit{stikord}}（禁整句）；Opgave 2 双人环节必须对\textbf{全部 4 张图}均明确表达偏好并用 \textit{fordi / derfor / synes} 简述理由，同时主动向搭档发问！
\end{enumerate}
\end{tipbox}

\subsubsection*{4.2 阅读 Opgave 1--4 逐题解析、干扰项拆解、全文中文精译与核心词表}

\begin{answerbox}
\textbf{\color{NCDNavy}【Opgave 1 逐空语法/语境解析、干扰词排雷与中文精译】}
\begin{itemize}
  \item \textbf{(1) hvor}：先行词为地点 \textit{en lille landsby på Fyn}，后接存在句 \textit{der kun bor...}（“在那里只住着大约四百人”）。此处从句副词 \textit{kun} 置于动词 \textit{bor} 之前（SAV 从句语序），直接排除需接主句语序的并列连词 \textit{men / så}。
  \item \textbf{(2) om}：固定搭配 \textit{usikker på, om...}（“不确定是否……”），引导间接是非疑问从句。词库中没有 \textit{at}，且语义表不确定性，故只能选 \textbf{om}。
  \item \textbf{(3) mens}：连接两个同时进行的动作，“一边听鸟儿歌唱，\textbf{一边}（\textit{mens}）一起吃早饭”。其他连词（\textit{fordi, hvis}）均无法构成合理的伴随时间关系。
  \item \textbf{(4) men}：语义急转折，“Signe 原以为乡下人各顾各的，\textbf{但是}（\textit{men}）第一个周末邻居就带着刚出炉的面包来敲门了”，且 \textit{allerede den første weekend bankede naboerne...} 为标准主句 V2 倒装，故排除从属连词 \textit{selvom / fordi}。
  \item \textbf{(5) til}：固定动词短语 \textit{vænne sig til noget}（“习惯于某事”）。\textit{under} 不与 \textit{vænne sig} 搭配。
  \item \textbf{(6) selvom}：让步转折，“尽管去办公室路程其实比以前更久，但他们依然很满意这个决定”。注意 \textit{turen til kontoret \underline{faktisk} tager...} 中 \textit{faktisk} 位于动词 \textit{tager} 之前（标准从句语序），既在语法上封死了 \textit{men}，又在逻辑上排除了表示因果的 \textit{fordi}（没人会“因为通勤变长而高兴”）。
  \item \textbf{6 个干扰词排雷}：\textit{fordi}（(4)(6) 逻辑均不成立）、\textit{hvis}（全文为已发生事实，无需假设条件）、\textit{ikke}（无法连接从句或搭配介词）、\textit{lige}（副词无合适空位）、\textit{så}（表“所以”，与 (4) 期望反转矛盾，且不能接 (6) 从句语序）、\textit{under}（与 \textit{vænne sig} 不搭配）。
\end{itemize}

\textbf{\color{NCDTeal}[Opgave 1 中文精译]}\\
在哥本哈根市中心一间昂贵的两居室公寓住了七年后，西涅和卡斯珀决定搬到乡下去。他们早就梦想能有更大的空间和一座花园，但在首都很难以合理的价格买到房子。春天的一个周末，他们去菲英岛一个小村庄看望朋友，那里只住着大约四百人。他们立刻被宁静的气氛和镇子周围绿油油的田野迷住了。尽管两人对漫长的通勤路程有些担心，他们还是买下了一栋院子里种着果树的旧砖房。第一个月西涅还有些拿不准自己是否会想念哥本哈根的咖啡馆和文化生活。夜里街道静悄悄的，早晨在花园里一边一起吃早饭，一边就能听见鸟儿歌唱。西涅原本以为乡下人会比较各顾各的，但才到第一个周末，邻居就捧着刚烤好的面包来敲门了。如今卡斯珀每周四天坐火车通勤，他也终于习惯了多出来的这段路程，因为他每周五都能居家办公。这对伴侣至今依然很庆幸自己的决定，尽管去办公室的路程其实比以前更费时间。乡村生活给他们的日常带来了更多的宁静。

\vspace{0.45em}
\hrule
\vspace{0.45em}
\textbf{\color{NCDNavy}【Opgave 2A \& 2B 逐空语篇指代解析与中文精译】}
\begin{itemize}
  \item \textbf{Opgave 2A (Mette \& Louise)}：\textbf{1. D}（Mette 说铁路施工火车停运 $\rightarrow$ Louise 感叹“真扫兴，那得坐铁路替代巴士，慢得要命” $\rightarrow$ 紧接 Mette 下一句 \textit{Det har du desværre ret i --- det tager næsten dobbelt så lang tid med togbus} 精准回扣）；\textbf{2. B}（Mette 提议改到周日 $\rightarrow$ Louise 说周日可以但傍晚要去看哥哥来不及既看电影又逛美术馆，提议“要不放弃美术馆？” $\rightarrow$ 下句 Mette 回答 \textit{Nej, så går vi jo glip af den nye udstilling...}）；\textbf{3. F}（Mette 提议上午看美术馆、放弃看电影 $\rightarrow$ Louise 同意并反问“你来订博物馆门票吗？” $\rightarrow$ 末句 Mette 回应“太好了，我来买 10:30 的博物馆票”）。\\
  \textbf{排雷}：\textbf{A}（说“喜欢替代巴士”与下句 \textit{Det har du desværre ret i} 矛盾）；\textbf{C}（周日全天上班与后续周日约定矛盾）；\textbf{E}（还停留在周六拿电影票，与改期周日去博物馆矛盾）。
  \item \textbf{Opgave 2B (Anders \& Camilla)}：\textbf{1. E}（Anders 问早上 9:30 的客户会议怎么办 $\rightarrow$ Camilla 说让 Martin 顶替并让 Anders 把资料发给\textbf{他} $\rightarrow$ 紧接下句 \textit{Jeg sender \underline{ham} straks alle noterne...} 精准回扣）；\textbf{2. A}（Camilla 提醒在考勤系统登记“儿童首日病假” $\rightarrow$ 下句 Anders 回答 \textit{Ja, jeg har allerede registreret det i vores tidsystem}）；\textbf{3. C}（Anders 问下午有没有特别优先事项 $\rightarrow$ Camilla 拜托他在下班前完成月度预算 \textit{månedsbudgettet} $\rightarrow$ 末句 Anders 回应 \textit{Det skal jeg nok gøre færdigt før kl. 16}）。\\
  \textbf{排雷}：\textbf{B}（说 Martin 在休假与下句发资料给“他”矛盾）；\textbf{D}（让 Anders 直接休息到周一与下句下午一点上 Teams 矛盾）；\textbf{F}（擅自给 Anders 妻子打电话与妻子中午回来接班的事实冲突）。
\end{itemize}

\textbf{\color{NCDTeal}[Opgave 2A \& 2B 中文精译]}\\
\textbf{2A（梅特与露易丝）}：梅特：“嗨露易丝！期待周六去奥胡斯看电影和逛美术馆吗？”｜露易丝(Z)：“超期待！我盼了一整周了。”｜梅特：“我也是，但我刚在 Rejseplanen 查到周六全天铁路施工，火车根本不开。”｜露易丝(1-D)：“唉真可惜！那我们估计得坐替代大巴了，那简直慢得像一个世纪。”｜梅特：“你说得太对了，坐替代大巴时间几乎翻倍。要不我们改成周日去？”｜露易丝(2-B)：“周日我倒是合适！但傍晚我得去看我哥，来不及同时看电影和逛展览了，要不把美术馆取消？”｜梅特：“不行，那样就错过新展览了，我可不想错过。那我们上午去美术馆、干脆放弃电影行吗？”｜露易丝(3-F)：“行啊，这计划听起来很棒！之后还能喝杯咖啡。你订博物馆门票吗？”｜梅特：“太好了，那我来买十点半的博物馆票，周日早上车站见！”\\
\textbf{2B（安德斯与卡米拉）}：安德斯：“早安卡米拉。我儿子诺亚早上发高烧了，我今天必须请‘孩子生病第一天’假。”｜卡米拉(Z)：“早安安德斯。真遗憾，祝他早日康复！你今天还能稍微居家办公吗？”｜安德斯：“可以，我太太中午回家接替照顾，所以下午我可以居家办公。不过我 9:30 的客户会议怎么办？”｜卡米拉(1-E)：“别担心，我问问马丁能不能替你开这个会。你能把资料发给他吗？”｜安德斯：“太感谢了，真是松了一口气！我马上把所有笔记和演示文稿发到他邮箱。”｜卡米拉(2-A)：“完美！记得在我们的系统里也登记一下孩子第一天病假，好吗？”｜安德斯：“好的，我已经在考勤系统里登记了，下午一点我会登录 Teams。下午有什么需要优先处理的事吗？”｜卡米拉(3-C)：“有的，如果你能在下班前把月度预算做完，那就帮大忙了。”｜安德斯：“我一定在四点前搞定。下午 Teams 见！”

\vspace{0.45em}
\hrule
\vspace{0.45em}
\textbf{\color{NCDNavy}【Opgave 3 逐段前瞻后顾逻辑与中文精译】}
\begin{itemize}
  \item \textbf{1. Svar: A}：空前说终于拿到 8 周会计事务所实习，空后用 \textbf{\textit{For selvom...}} 解释“因为虽然她在本国做过办公室工作，但所有丹麦法规和词汇对她来说都是全新的”，故空位必须同时表达“既开心又有一点紧张”的矛盾心理（A）。B/C/D 均与拿到实习或语言陌生矛盾。
  \item \textbf{2. Svar: B}：空前介绍头几天要看旧发票、学软件，空后用 \textit{Heldigvis siger Mette, at hun altid må spørge, hvis der er et fagord...} 安慰她，说明空位必须交待“由于专业词多，比她预想的费时间”（B）。A 说“很简单”与后文安慰逻辑冲突。
  \item \textbf{3. Svar: C}：空前讲八位同事午休说话很快并开内部玩笑，空后用 \textbf{\textit{Derfor beslutter hun sig for...}} 引出她决定主动坐到友善同事身边发问，故空位须是导致她“不知何时插话”的心理原因（C）。A“喜欢独自吃饭”与后句坐到同事旁冲突。
  \item \textbf{4. Svar: C}：空前发现文件夹里少了一大笔钱，空后 Mette 称赞她细心发现了错误，说明中间缺失的正是“找了半小时终于找到丢失单据并拿给 Mette 看”（C）。A/B 隐瞒不报则不可能得到 Mette 表扬。
  \item \textbf{5. Svar: B}：空前说被叫进老板办公室、心情忐忑，空后说她“带着大大的笑容回家，觉得稳定工作的梦想真的有望实现”，故老板宣布的必然是“团队很满意，下个月给她一份兼职代班合同”（B）。
\end{itemize}

\textbf{\color{NCDTeal}[Opgave 3 中文精译]}\\
(1) 法蒂玛在欧登塞住了两年，一直梦想在丹麦从事会计和财务工作。上个月她通过了丹麦语考试，终于在当地一家会计师事务所拿到了为期八周的实习岗位。第一天上班前，她既非常开心又有点紧张。因为虽然她在家乡做过办公室工作，但所有丹麦的规章制度和专业词汇对她来说都是全新的。(2) 周一早上，她的联系人梅特热情迎接了她，带她参观办公室并把她介绍给其他员工。头几天法蒂玛需要阅读旧发票并熟悉公司的数字系统。因为很多专业词汇很难，这件事比她预想的要花更多时间。幸好梅特说，只要遇到不懂的术语或任务，她随时都可以问。(3) 两周后法蒂玛更熟练地掌握了电脑软件，但食堂里的午休时间依然让人觉得有点吃力。八位同事围坐吃饭时语速很快，还会笑一些圈内梗。这有时让她拿不准自己什么时候该开口插话。因此她决定坐到一位友善的同事旁边，主动问几个关于周末计划的问题。(4) 实习期过半时，法蒂玛获准独立核对当月收据并把数字录入账目。一天下午她突然发现其中一个文件夹里少了一大笔金额。找了半小时后，她幸好找到了丢失的凭证并拿给梅特看。梅特对法蒂玛发现这个错误赞不绝口，并表扬她非常细心。(5) 八周快结束时，法蒂玛被叫进老板办公室谈话。她很紧张老板会对她实习期间的表现作何评价。老板告诉她，因为团队非常喜欢她，决定从下个月起给她提供一份兼职替班合同。法蒂玛面带大大的笑容回到家，因为现在她知道，在丹麦找到稳定工作的梦想真的可以实现。

\vspace{0.45em}
\hrule
\vspace{0.45em}
\textbf{\color{NCDNavy}【Opgave 4 逐题同义定位与三人全文中文精译】}
\begin{itemize}
  \item \textbf{1. Svar: B (Karim)}：\textit{professionel hjælp udefra til rengøringen} $\leftrightarrow$ Karim：\textit{betale et rengøringsfirma, der kommer hver anden fredag}。
  \item \textbf{2. Svar: A (Stine)}：\textit{en fast skriftlig plan i køkkenet} $\leftrightarrow$ Stine：\textit{et fast ugeskema på køleskabet}（冰箱上的固定周计划表）。
  \item \textbf{3. Svar: C (Elena)}：\textit{overtaget al madlavningen, fordi partneren ikke var god til det} $\leftrightarrow$ Elena：Frederik 两次把菜烧糊且讨厌看食谱，现在由 Elena \textit{står for al madlavningen}。
  \item \textbf{4. Svar: B (Karim)}：\textit{bede sine børn flere gange} $\leftrightarrow$ Karim：\textit{minde dem om det flere gange}（不得不多次提醒三个青少年）。Stine 的孩子虽有固定家务但未提需反复催促。
  \item \textbf{5. Svar: C (Elena)}：\textit{irriterende, at hjemmearbejde fører til ekstra husarbejde} $\leftrightarrow$ Elena：在家办公三天时常顺手洗衣服，\textit{Det irriterer mig faktisk lidt, for hjemmearbejde er jo stadig rigtigt arbejde!}
  \item \textbf{6. Svar: A (Stine)}：\textit{køber ind..., mens partneren laver det meste af aftensmaden} $\leftrightarrow$ Stine：她负责洗衣和采购（\textit{indkøbene}），而 Nikolaj 每周做五天晚饭。Karim 自己既买菜又做工作日全套晚饭，不符“伴侣做大部分晚饭”。
  \item \textbf{7. Svar: B (Karim)}：\textit{overtog madlavningen i hverdagene, fordi partneren fik skiftende arbejdstider} $\leftrightarrow$ Karim：妻子 Sarah 当护士轮值夜班（\textit{skiftende nattevagter}）后，他开始负责工作日买菜和做热晚饭。
\end{itemize}

\textbf{\color{NCDTeal}[Opgave 4 中文精译]}\\
\textbf{A. 斯蒂娜}：以前两个孩子还很小的时候我做兼职，几乎包揽了所有家务。可那导致我和丈夫尼古拉晚上经常吵架，因为我总是累坏了。现在我们俩都全职工作，因此在冰箱上贴了一份固定的周计划表。尼古拉负责吸尘和打扫卫生间，我则负责洗衣服和采购。作为平衡，尼古拉特别喜欢下厨，所以他每周做五天晚饭，我只在周末简单做点吃的。我们六岁和九岁的孩子也有固定家务：必须自己收拾房间，并每周帮忙清空几次洗碗机。我觉得家庭平权不是去斤斤计较每一分钟，而是让父母双方晚上都能有时间喘口气放松一下。\\
\textbf{B. 卡里姆}：在我童年时都是妈妈包办所有做饭和打扫，所以刚和萨拉结婚时，我老实说也以为家里会是这样。但萨拉开始做护士、要倒班上夜班后，情况很快变了。如今几乎总是由我买菜并在每个工作日做热晚饭，因为她下班后常常精疲力竭。我们俩都不想把宝贵的周末花在擦地板上，所以选择花钱请保洁公司隔周五来打扫一次。我们的三个青少年孩子不爱做饭，不过他们轮流割草坪、倒垃圾和叠衣服——虽然我还是得提醒他们好几次。在我看来，男人在家庭里应当承担和女人一样多的责任。\\
\textbf{C. 埃琳娜}：我和丹麦男友弗雷德里克住在一起，还没有孩子。因为我每周有三天居家办公，午休时不知不觉就会去洗一桶衣服，而弗雷德里克每天要通勤两小时。这其实让我有点烦，因为居家办公也是正经上班啊！以前我们试过轮流做饭，但弗雷德里克两次把菜烧糊，而且讨厌照着食谱做。现在所有做饭都归我，因为我本来就喜欢下厨。作为交换，弗雷德里克事后总是负责洗碗、擦所有窗户，并把垃圾送去回收站。我认为好的平权是按各自最擅长的事情来分工，而不是两个人非得干一模一样的活儿。将来有了孩子，也要从小教他们搭把手。
\end{answerbox}

\subsubsection*{本卷核心生词精表（Gloser --- Modultest 4 · 全 15 词 B1 精编）}

\begin{glosertable}
名词 & en landsby & \textit{(landsbyen, landsbyer, landsbyerne)} & 村庄，乡村小镇 \\
名词 & et foreningsliv & \textit{(foreningslivet, ingen flertal)} & 社团生活，协会活动 \\
动词 & at pendle & \textit{(pendler, pendlede, har pendlet)} & 跨城通勤，往返上下班 \\
名词 & en afgang & \textit{(afgangen, afgange, afgangene)} & 发车班次，出发时间 \\
名词 & en forsinkelse & \textit{(forsinkelsen, forsinkelser, forsinkelserne)} & 延误，晚点 \\
名词 & en erstatning & \textit{(erstatningen, erstatninger, erstatningerne)} & 赔偿，补偿金 \\
名词 & en refusion & \textit{(refusionen, refusioner, refusionerne)} & 退款，退费 \\
名词 & en praktikplads & \textit{(praktikpladsen, praktikpladser, praktikpladserne)} & 实习岗位 \\
名词 & et regnskab & \textit{(regnskabet, regnskaber, regnskaberne)} & 账目，财务报表 \\
名词 & en arbejdsdeling & \textit{(arbejdsdelingen, arbejdsdelinger, arbejdsdelingerne)} & 劳动分工，家务分配 \\
名词 & ligestilling & \textit{(ligestillingen, ingen flertal)} & 平等，平权（常指性别平等） \\
名词 & en pligt & \textit{(pligten, pligter, pligterne)} & 义务，分内家务/职责 \\
形容词 & omhyggelig & \textit{(omhyggeligt, omhyggelige)} & 细心的，认真仔细的 \\
动词 & at tilpasse sig & \textit{(tilpasser sig, tilpassede sig, har tilpasset sig)} & 适应，主动调整适应 \\
名词 & en ferieform & \textit{(ferieformen, ferieformer, ferieformerne)} & 度假方式，旅游形式 \\
\end{glosertable}

\subsubsection*{4.3 写作双任务双档范文（Opgave 1 投诉信 \& Opgave 2 私人回信）}

\begin{answerbox}
\textbf{\color{NCDNavy}【Opgave 1 范文一：B1 稳过版（105 ord · 四要点清晰齐全）】}\\
\textit{Til: Kundeservice, NordBus \quad|\quad Emne: Klage over busrejse fra København til Aalborg}

\vspace{0.25em}
\textit{Kære NordBus}

\vspace{0.2em}
\textit{Jeg skriver til jer for at klage over en meget dårlig busrejse i går den 18. november. Jeg rejste med jeres bus fra København til Aalborg klokken 09.15.}

\vspace{0.2em}
\textit{For det første var bussen næsten to timer forsinket, og vi fik desværre ingen besked. Derudover var mit sæde ødelagt, og toilettet i bussen virkede ikke under turen.}

\vspace{0.2em}
\textit{Da jeg talte med chaufføren i pausen, sagde han, at han ikke kunne gøre noget. Bagefter ringede jeg til jeres kundeservice, men de ville heller ikke hjælpe mig.}

\vspace{0.2em}
\textit{Derfor håber jeg, at I vil give mig en undskyldning og betale billetten tilbage.}

\vspace{0.2em}
\textit{Med venlig hilsen\\
Sofie Jensen}

\vspace{0.25em}
\small\textbf{点评（正文 105 词）}：句式稳健零硬伤，精准兑现全部 4 个要点（乘车时间与起终点 / 延误无通知 + 坏座椅 + 故障厕所 / 司机与客服推诿反应 / 道歉与全额退票诉求），全程对客运公司统一使用 \textit{I / jer / jeres}，正确运用从句否定前置（\textit{at han ikke kunne}）。

\vspace{0.45em}
\hrule
\vspace{0.45em}

\textbf{\color{NCDRed}【Opgave 1 范文二：B1+ 高分示范版（137 ord · 亮点句法全调度）】}\\
\textit{Til: Kundeservice, NordBus \quad|\quad Emne: Klage over forsinkelse og dårlig service}

\vspace{0.25em}
\textit{Kære NordBus}

\vspace{0.2em}
\textit{Jeg vil hermed klage over en meget utilfredsstillende rejse mellem København og Aalborg i går den 18. november med afgang kl. 09.15.}

\vspace{0.2em}
\textit{Allerede fra starten opstod der problemer. Bussen \da{blev} nemlig \da{forsinket} med næsten to timer, uden at passagererne \da{blev informeret} undervejs. Da vi endelig kørte, opdagede jeg, at mit sæde var knækket, og at toilettet var låst på hele den lange tur.}

\vspace{0.2em}
\textit{Da jeg gjorde chaufføren opmærksom på fejlene, reagerede han desværre uhøfligt og bad mig om at sætte mig igen. Bagefter kontaktede jeg jeres kundeservice, men der \da{blev ikke taget} ansvar for situationen, \dabold{hvilket} jeg er meget skuffet over.}

\vspace{0.2em}
\textit{\dabold{Har man betalt} fuld pris for en behagelig rejse, \dabold{forventer man} ordentlige forhold. Jeg håber derfor, at I vil give mig en fuld refusion af billetten samt et gavekort.}

\vspace{0.2em}
\textit{Med venlig hilsen\\
Sofie Jensen}

\vspace{0.25em}
\small\textbf{高分点评（正文 137 词）}：全面调度 B1 高阶句法：\ding{172} 连续精准使用 \textit{blive} 被动（\textit{blev forsinket, blev informeret, blev ikke taget ansvar}）客观陈述具体投诉事件；\ding{173} \textit{uden at} 否定伴随从句与 \textit{bad mig om at} 介词搭配；\ding{174} 整句关系从句 \textit{..., hvilket jeg er meget skuffet over} 提升逻辑连贯度；\ding{175} 末段以“无 \textit{hvis} 条件倒装”（\textit{Har man betalt..., forventer man...}）收尾索赔，正式得体且说服力极强。

\vspace{0.45em}
\hrule
\vspace{0.45em}

\textbf{\color{NCDNavy}【Opgave 2 范文一：B1 稳过版（111 ord · 远超 90 词底线）】}\\
\textit{Til: Line \quad|\quad Emne: Svar: Livet i Danmark nu}

\vspace{0.25em}
\textit{Kære Line}

\vspace{0.2em}
\textit{Tak for din mail! Det er dejligt at høre fra dig.}

\vspace{0.2em}
\textit{Det, der overrasker mig mest i Danmark, er, hvor meget folk stoler på hinanden, og at man bruger fornavn til sin chef.}

\vspace{0.2em}
\textit{Heldigvis er mange ting blevet nemmere nu. Jeg forstår endelig MitID og e-Boks, og jeg tør tale dansk, når jeg handler ind eller snakker med naboerne.}

\vspace{0.2em}
\textit{Til gengæld synes jeg stadig, at den hurtige udtale er svær, når flere danskere taler sammen. Det er også svært at få danske venner uden for skolen.}

\vspace{0.2em}
\textit{I de næste måneder skal jeg til modultest i dansk, og bagefter vil jeg søge en praktikplads på et kontor.}

\vspace{0.2em}
\textit{Kærlig hilsen\\
Maya}

\vspace{0.25em}
\small\textbf{点评（正文 111 词）}：自然呼应 Line 的全部 4 个问题（最惊讶之处 / 变容易的事 / 仍觉得困难之处 / 未来数月计划），字数稳超 90 词要求（不含抬头标题已达 111 词），称呼与落款温馨得体。

\vspace{0.45em}
\hrule
\vspace{0.45em}

\textbf{\color{NCDRed}【Opgave 2 范文二：B1+ 高分示范版（132 ord · 文化观察与句法示范）】}\\
\textit{Til: Line \quad|\quad Emne: Svar: Livet i Danmark nu}

\vspace{0.25em}
\textit{Kære Line}

\vspace{0.2em}
\textit{Hvor er det hyggeligt at høre fra dig! Mit liv i Danmark føles meget anderledes nu end i begyndelsen.}

\vspace{0.2em}
\textit{\da{Det, der} stadig overrasker mig mest her, er den store tillid og det flade hierarki. Man bruger fornavn til chefen, og forældre lader trygt barnevogne stå ude på gaden.}

\vspace{0.2em}
\textit{Samtidig er hverdagen blevet væsentlig nemmere. Jeg behøver ikke længere at oversætte hver besked i e-Boks, og fordi jeg læser danske nyheder, føler jeg mig langt mere hjemme.}

\vspace{0.2em}
\textit{Alligevel synes jeg, at det sociale liv er svært. \dabold{Til trods for at} danskerne er venlige på jobbet, \da{planlægges} næsten alle aftaler uger i forvejen, og de bløde d'er driller mig stadigvæk!}

\vspace{0.2em}
\textit{I de næste måneder skal jeg bestå Modul 3, og derefter \da{regner jeg med, at} jeg skal i virksomhedspraktik.}

\vspace{0.2em}
\textit{Kærlig hilsen\\
Maya}

\vspace{0.25em}
\small\textbf{高分点评（正文 132 词）}：紧扣 NCD 3 文化深读亮点：\ding{172} 精准点出 \textit{tillid} 与 \textit{det flade hierarki}；\ding{173} 运用习惯性 \textit{-s} 被动 \textit{planlægges} 对比丹麦人的社交预约习惯；\ding{174} 亮出高级让步连词 \textit{Til trods for at} 与“介词 + \textit{at}-从句”（\textit{regner jeg med, at...}），自然流畅地达标 B1+ 水准。
\end{answerbox}

\subsubsection*{4.4 口语示范稿（Mindmap 4A \& 4B 约 150 词陈述 + 语音 + Del 2 问答与 Opgave 2 四图互动对话）}

\begin{answerbox}
\textbf{\color{NCDNavy}【Mindmap 4A：Arbejdsdeling i min familie（Del 1 约 150 词脱稿陈述示范 + 核心音标）】}

\vspace{0.18em}
\textit{Jeg vil gerne fortælle om arbejdsdelingen \ipa{ˈɑːbɑjdsˌdeːleŋən} i min familie. Jeg bor i en lejlighed sammen med min partner Thomas og vores datter på otte år. Da vi begge arbejder på fuldtid, er vi nødt til at samarbejde i hverdagen.}

\vspace{0.18em}
\textit{Med hensyn til madlavning og indkøb \ipa{ˈenˌkʰøːˀb} skiftes vi til at lave aftensmad hver anden dag. Hver søndag laver vi en madplan, og om lørdagen handler vi sammen. Rengøringen klarer vi med en fast lørdagsrutine: Thomas støvsuger og gør badeværelset rent, mens jeg vasker og lægger tøj sammen.}

\vspace{0.18em}
\textit{Vores datter har også faste pligter. Hun skal rydde op på sit værelse, dække bord og tømme opvaskemaskinen \ipa{ˈʌbvasgəmaˌsgiːnən}, før hun ser fjernsyn. Økonomien \ipa{økonoˈmiːˀən} styrer vi via en fælles budgetkonto til faste udgifter, og vi deler en digital familiekalender på mobilen.}

\vspace{0.18em}
\textit{Jeg synes, ligestilling \ipa{ˈliːəˌsdeleŋ} i hjemmet er meget vigtigt. Det handler om, at begge voksne har fritid, og at vi giver vores datter et godt forbillede.}

\vspace{0.2em}
\small\textbf{[参考译文]}：我想谈谈我家的家务分工。我和伴侣托马斯以及我们八岁的女儿住在一套公寓里。因为我们俩都全职工作，平日里必须互相配合。在做饭和采购方面，我们隔天轮流做晚饭；每周日一起制定食谱，周六一起去采购。打扫卫生我们也有固定的周六分工流程：托马斯负责吸尘和清理卫生间，我负责洗衣服和叠衣服。我们的女儿也有固定任务：看电视前必须整理房间、摆好餐桌并清空洗碗机。财务方面我们用共同的预算账户管理固定开支，并在手机上共享数字家庭日历。我认为家庭平权非常重要，这意味着两位成年人都拥有空闲时间，同时也能给女儿树立好榜样。

\vspace{0.25em}
\textbf{\color{NCDRed}Del 2 考官追问与示范回答：}
\begin{enumerate}
  \item \textbf{Q:} \textit{Du siger, at I skiftes til at lave mad hver anden dag. Hvad laver du helst, når det er din tur?}\\
  \textbf{A:} \textit{Når det er min tur, laver jeg helst nemme, varme retter som kylling i karry eller pastaretter med masser af grøntsager, fordi vi alle tre kommer trætte hjem.}
  \item \textbf{Q:} \textit{Hvordan reagerer jeres datter på sine pligter --- gør hun dem altid uden problemer?}\\
  \textbf{A:} \textit{Ikke altid! Nogle gange må vi minde hende om opvaskemaskinen et par gange, men som regel gør hun det gerne, især når vi sætter glad musik på.}
  \item \textbf{Q:} \textit{Er der nogle huslige opgaver, som du eller din partner slet ikke kan lide at lave?}\\
  \textbf{A:} \textit{Ja, ingen af os kan lide at pudse vinduer eller stryge skjorter, så vi stryger faktisk næsten aldrig tøj i hverdagen.}
  \item \textbf{Q:} \textit{Hvordan har arbejdsdelingen i dit hjem ændret sig i forhold til, hvordan du selv voksede op?}\\
  \textbf{A:} \textit{Da jeg var barn, lavede min mor næsten alt i køkkenet alene. I dag er det heldigvis meget mere naturligt, at mænd og kvinder deler både ansvaret og opgaverne ligeligt.}
\end{enumerate}

\vspace{0.4em}
\hrule
\vspace{0.4em}

\textbf{\color{NCDNavy}【Mindmap 4B：Mine erfaringer med at lære dansk（Del 1 约 150 词脱稿陈述示范 + 核心音标）】}

\vspace{0.18em}
\textit{Jeg vil fortælle om mine erfaringer \ipa{ɐˈfɑːˀeŋɐ} med at lære dansk. Jeg går på et aftenhold på sprogskolen to gange om ugen på DU3, Modul 3, hvor vi er seksten kursister i klassen fra forskellige lande.}

\vspace{0.18em}
\textit{Ud over undervisningen bruger jeg cirka femogfyrre minutter dagligt på lektier. Jeg lytter til danske podcasts i toget, læser korte danske nyheder online og skriver nye gloser ned i en notesbog. Det sværeste ved sproget er helt klart udtalen \ipa{ˈuðˌtˢɑːlən} --- især de bløde d'er og vokalerne. I grammatikken øver jeg meget på ordstilling \ipa{ˈoɐ̯ˌsdeleŋ} i ledsætninger og passivformer.}

\vspace{0.18em}
\textit{I hverdagen bruger jeg dansk så meget som muligt, når jeg handler i butikker eller snakker med kollegaer i frokostpausen. Min største udfordring \ipa{ˈuðˌfɒːdʁeŋ} er, når folk taler hurtigt, fordi jeg nogle gange bliver genert og skifter til engelsk. Alligevel er jeg stolt af min udvikling \ipa{ˈuðˌveɡleŋ}, for jeg føler mig mere selvsikker nu og håber at nå PD3.}

\vspace{0.2em}
\small\textbf{[参考译文]}：我想分享一下自己学习丹麦语的经历。我在语言学校每周上两次夜校班，目前在 DU3 Modul 3，班里有十六位来自不同国家的同学。除了课堂学习，我每天大约花四十五分钟做作业。我会在火车上听丹麦语播客、在网上读简短的丹麦语新闻，并把新单词记在笔记本里。这门语言最难的部分无疑是发音——尤其是软 d 和元音。语法上我重点练习从句语序和被动语态。日常生活中，我在商店购物或午休和同事聊天时尽量多用丹麦语。最大的挑战是别人语速快的时候，因为有时我会害羞并切回英语。尽管如此，我对自己的进步感到自豪，现在自信多了，希望能顺利达到 PD3。

\vspace{0.25em}
\textbf{\color{NCDRed}Del 2 考官追问与示范回答：}
\begin{enumerate}
  \item \textbf{Q:} \textit{Du nævner, at du lytter til podcasts i toget. Hvilke podcasts eller medier kan du anbefale?}\\
  \textbf{A:} \textit{Jeg kan varmt anbefale DR's nyheder på let dansk samt podcasten ”Dansk i ørerne”, fordi værten taler tydeligt og forklarer spændende ting om dansk kultur.}
  \item \textbf{Q:} \textit{Hvad gør du, når en dansker pludselig skifter til engelsk midt i en samtale?}\\
  \textbf{A:} \textit{Før i tiden skiftede jeg bare med, men nu smiler jeg og siger pænt: ”Jeg vil faktisk rigtig gerne øve mit dansk, så lad os fortsætte på dansk!” Det respekterer folk altid.}
  \item \textbf{Q:} \textit{Hvad synes du har været det vigtigste, I har lært på Modul 3 indtil nu?}\\
  \textbf{A:} \textit{Det vigtigste for mig har været at kunne skrive formelle klager og fortælle sammenhængende om samfundsemner i stedet for bare at bruge korte hovedsætninger.}
  \item \textbf{Q:} \textit{Hvordan vil du fortsætte med at bruge og forbedre dit danske efter modultesten?}\\
  \textbf{A:} \textit{Jeg vil gerne starte som frivillig i en lokal idrætsklub eller læsekreds, så jeg får endnu flere danske bekendte at snakke med i min fritid.}
\end{enumerate}
\end{answerbox}

\vspace{0.35em}
\textbf{\color{NCDNavy}\large 【Mundtlig Opgave 2：双人四图互动示范逐字稿（EMNE: Ferie og rejser）】}

\dialogturn{1}{Kursist A}{Skal vi starte med billede et? Jeg elsker virkelig at holde ferie i et sommerhus ved kysten, fordi der er så roligt, og man kan hygge sig med familien uden stress. Hvad synes du om sommerhusferie?}{ˈsʌmɐˌhuːˀs / ˈkʰøsdən / ˈhyɡə sɑj}{我们先从图一开始好吗？我真的很喜欢去海边度假屋度假，因为那里很安静，可以毫无压力地和家人共度温馨时光。你觉得度假屋旅行怎么样？}

\dialogturn{2}{Kursist B}{Jeg kan også rigtig godt lide sommerhuse, især om sommeren, når man kan gå ned til stranden og grille om aftenen. Men hvis det regner hele ugen, kan det godt blive lidt kedeligt. Hvad med billede to --- kan du lide at tage på storbyferie?}{ˈsdoɐ̯ˌbyːˌfeɐ̯jə / ˈʁɑjnɐ}{我也很喜欢度假屋，特别是夏天可以去海边散步、晚上烧烤。不过如果连下一整周雨，可能会有点无聊。那图二呢——你喜欢大城市度假吗？}

\dialogturn{3}{Kursist A}{Ja, absolut! Jeg foretrækker faktisk storbyferie i Europa, fordi jeg interesserer mig meget for kunst, museer og god mad. Selvom det kan være lidt dyrt og larmende, oplever man altid noget nyt. Foretrækker du storbyferie eller naturferie som på billede tre?}{foɐˈtˢʁɛɡɐ / muˈseːɐ / nɑˈtˢuːɐ̯ˌfeɐ̯jə}{是的，非常喜欢！其实我更偏爱欧洲大城市旅行，因为我对艺术、博物馆和美食很感兴趣。虽然可能有点贵又有点吵，但总能体验到新鲜事物。你更喜欢大城市度假还是像图三那样的自然度假？}

\dialogturn{4}{Kursist B}{Jeg er faktisk helt vild med aktiv naturferie som på billede tre! Når man vandrer i bjergene eller sover i telt i skoven, får man masser af frisk luft og glemmer arbejdet fuldstændig. Det kræver selvfølgelig godt udstyr, og man kan blive øm i benene, men jeg synes, det er det hele værd. Har du nogensinde prøvet camping?}{ˈvɑndʁɐ / ˈuðˌsdyːɐ̯ˀ / ˈømˀ}{其实我超爱图三这种户外自然度假！在山里徒步或在森林里露营时，能呼吸大量新鲜空气，把工作彻底抛在脑后。当然这需要好装备，腿也可能会酸，但我觉得完全值得。你试过露营吗？}

\dialogturn{5}{Kursist A}{Kun én gang, og ærligt talt savnede jeg en blød seng og et varmt bad! Hvad synes du egentlig om det sidste billede med staycation --- altså at holde ferie derhjemme og tage på korte dagsture?}{ˈæɐ̯lid ˈtˢɑːˀld / ˈdɑwsˌtˢuːɐ}{只试过一次，说实话我很想念软床和热水澡！那你对最后一张图——居家度假，也就是待在家里然后去短途一日游——怎么看？}

\dialogturn{6}{Kursist B}{Jeg synes faktisk, det er en ret smart og klimavenlig løsning, hvis man ikke har så mange penge eller bare trænger til at sove længe. Man sparer meget transporttid, men risikoen er, at man bare kommer til at ordne vasketøj derhjemme. Hvad tænker du om at holde ferie hjemme?}{ˈkʰliːmaˌvɛnli / ˈvɑsgəˌtˢʌj}{其实我觉得如果预算不多或者只是想睡个懒觉，这是个既聪明又环保的选择。能省下很多路途时间，但风险是一不小心就在家顺手洗起衣服做起家务了。你觉得在家度假怎么样？}

\dialogturn{7}{Kursist A}{Det er jeg helt enig i! Hvis jeg bliver hjemme hele ferien, føles det ikke rigtig som en pause. Derfor vil jeg helst rejse et par dage ud af byen --- selv bare til en anden landsdel i Danmark.}{ˈeːni / ˈlɑnsˌdeːˀl}{我完全同意！如果整个假期都待在家里，感觉不太像真正放假。所以我还是更愿意离开城市出去玩几天——哪怕只是去丹麦的另一个地区也好。}

\begin{answerbox}
\textbf{\color{NCDRed}【Opgave 2 --- Del 2 个人深入三问示范回答（Individuel opgave）】}
\begin{enumerate}
  \item \textbf{Q:} \textit{Hvordan kan du bedst lide at holde ferie --- og hvorfor?}\\
  \textbf{A:} \textit{Jeg kan bedst lide en blanding af storbyferie og lidt afslapning ved havet. Jeg elsker at se seværdigheder og prøve lokale retter, men efter et par travle dage har jeg også brug for at sidde stille med en god bog.}
  \item \textbf{Q:} \textit{Hvad er vigtigst for dig på en ferie (fx pris, vejr, afslapning eller nye oplevelser)?}\\
  \textbf{A:} \textit{For mig er nye oplevelser og godt selskab allervigtigst. Vejret kan man jo ikke styre --- især ikke i Danmark --- men hvis man rejser sammen med gode venner eller familie og har et rimeligt budget, får man næsten altid en god tur.}
  \item \textbf{Q:} \textit{Kan du fortælle lidt om den bedste ferie, du har haft for nylig? (hvor, hvornår, med hvem, hvad lavede I)}\\
  \textbf{A:} \textit{Sidste efterår tog min mand og jeg fire dage til Lissabon i Portugal. Vi havde fantastisk solskin på tyve grader i oktober, og vi gik rundt i de gamle gader med sporvogn, lyttede til fado-musik og spiste frisk fisk ved havnen. Det bedste var, at vi slet ikke havde travlt og bare nød stemningen sammen.}
\end{enumerate}
\end{answerbox}
